\documentclass[final,a4paper,12pt]{article}

\usepackage{notes}
\usepackage{placeins}

\title{Geophysics III}
\author{Problem Set \# 5}
\date{}

\begin{document}
\maketitle

The gist of the tectonic evolution problem is to develop a simple (back of the
envelope) model of otherwise relatively complex geologic processes.  This short
document provides a simplified description of diking and the type of detail I am
looking for in an appropriate answer.
\begin{figure}[hbt]
    \centering
    \includegraphics[]{figures/tectonic_processes_arc.eps}
    \caption{Hypothetical tectonic evolution of an arc and rift.}
    \label{fig:arc}
\end{figure}

Consider extension of the lithosphere by diking (Figure~\ref{fig:arc}).  The
process of adding dikes to the lithosphere transports significant quantities of
heat to the surface.  While the temperatures of the melts are similiar to those
found within the mantle over a sufficiently large region, the dikes raise
temperatures near the surface.  The greater the percentage of lithosphere
expanded by dikes, the the higher the average temperatures.  Dikes raise
temperatures throughout the lithosphere as well, even if they are only emplaced
near the surface.  We will assume the initial geotherm, for the sake of
simplicity, is nearly linear between a new higher temperature very near the
surface and the asthenospheric temperature at the base of the lithosphere
(Figure~\ref{fig:diking}).  The curvature comes from the influence of heat
production.
\begin{figure}[hbt]
    \centering
    \includegraphics[]{figures/diking.eps}
    \caption{A simple model of diking and its thermal evolutions.}
    \label{fig:diking}
\end{figure}

As for the temporal evolution, temperatures rapidly cool near the surface to a
fixed temperature $T_0$.  As a result the heat flow rapidly falls, much in the
same way the estimated heat flow of a rift would fall.  At the base of the
lithosphere temperature changes very little.  As cooling progresses, the
temperature begins to change more and more slowly as the cooling geotherm
approaches the equilibrium geotherm.  At this point, heat flow begins to
approach an asymptotic heat flow defined by the equilibrium geotherm.
Timescales for cooling are proportional to the square of the lithospheric
thickness, i.e., $\tau \propto L^2$.

The cooling of a rift and arc will be similar, but somewhat different than
diking, while the thermal evolution of a thrust is more complicated.  Draw a
picture of the tectonic evolution of these terranes as well to help aid your
description.

\end{document}
